Na seção anterior vimos que, para identificar um efeito causal, precisamos comparar o que de fato ocorreu com o que teria acontecido na ausência da política. Chamamos esse cenário não observado de \emph{contrafactual}. Em um experimento aleatorizado, essa comparação é simples: o sorteio garante que os grupos sejam equivalentes em média, de modo que qualquer diferença de resultado pode ser atribuída ao tratamento. Mas, fora do ambiente experimental, essa equivalência precisa ser \textbf{reconstruída a partir dos dados}.

A regressão é a ferramenta mais usada para esse propósito. No entanto, como discutido na seção anterior, a validade causal dessa abordagem depende de hipóteses fortes que muitas vezes não são críveis em dados observacionais. Por esse motivo, é útil considerar estratégias alternativas para lidar com o viés de seleção.

O pareamento nasce exatamente desse incômodo. Em vez de depender de uma forma funcional ou de um ajuste paramétrico, ele parte de uma ideia mais direta: \textbf{se quisermos saber o que teria acontecido com os tratados, comparemos com não tratados que se pareçam com eles}. A lógica é simples, mas poderosa. Se cada beneficiário do programa tiver uma ``dupla'' no grupo de controle com características muito próximas (renda, idade, escolaridade, região, etc.), então a diferença média entre os dois grupos fornece uma boa aproximação do efeito causal.

A diferença central em relação à regressão está na ordem lógica do procedimento. No pareamento, a comparabilidade entre tratados e controles é construída primeiro, por meio da restrição explícita do conjunto de comparação; a estimação do efeito vem depois, como uma simples diferença de médias. Na regressão, essas duas etapas --- construção da comparabilidade e estimação do efeito --- ocorrem simultaneamente, sendo mediadas pela forma funcional imposta ao modelo.


Podemos pensar o pareamento como uma tentativa de “recriar” o sorteio de um experimento, mas usando as semelhanças observadas entre as unidades. Ele busca equilibrar os grupos \emph{antes} da estimação, de modo que, ao final, o único fator que realmente varie entre eles seja o tratamento. O avaliador não está modelando a relação entre variáveis, e sim \textbf{selecionando comparações justas}.

\section{A ideia central do pareamento}

Na Parte 2 deste guia, mostramos que, em dados não-experimentais, a simples comparação entre tratados e não tratados não identifica, em geral, um efeito causal. O problema fundamental é que a decisão de participar do tratamento não é aleatória: ela depende de características observáveis e não observáveis dos indivíduos. Como consequência, em geral,
$$
\mathbb{E}[Y_i(0) \mid D_i = 1] \neq \mathbb{E}[Y_i(0) \mid D_i = 0],
$$

o que invalida a interpretação causal da diferença de médias.

A resposta conceitual a esse problema foi apresentada na Seção 2.2 por meio do modelo de resultados potenciais: se, ao condicionarmos em um conjunto adequado de variáveis  $X_i$, a participação no tratamento passar a ser independente dos resultados potenciais, então a identificação causal se torna possível. No exemplo do NSW, entram nesse conjunto escolaridade, idade, histórico de renda, desemprego prévio, raça, entre outras. Essa é a Hipótese de Independência Condicional (\textit{Conditional Independence Assumption} - CIA):
$$
(Y_i(1),Y_i(0)) \perp D_i \mid X_i
$$

A ideia central do pareamento nasce exatamente dessa observação: se conseguirmos tornar tratados e controles comparáveis em termos das covariadas relevantes $X_i$, então a comparação entre eles volta a ter interpretação causal. Assim, o pareamento não deve ser entendido, em primeiro lugar, como uma técnica computacional, mas como uma estratégia de identificação: ele busca recriar, dentro dos dados observacionais, a situação de comparabilidade que um experimento aleatorizado geraria automaticamente.

Sob a hipótese de independência condicional, o objeto causal fundamental deixa de ser diretamente o ATE ou o ATT agregados. O foco passa a ser os efeitos causais condicionais, definidos para cada perfil observável $X_i = x$:
\[
\text{ATE}(x) := \mathbb{E}\!\left[ Y_i(1) - Y_i(0) \mid X_i = x \right], 
\quad 
\text{ATT}(x) := \mathbb{E}\!\left[ Y_i(1) - Y_i(0) \mid D_i = 1, X_i = x \right].
\]

O pareamento não deve ser interpretado apenas como uma técnica alternativa de estimação, mas como uma abordagem que prioriza a comparabilidade entre tratados e controles no espaço das covariadas antes da comparação de resultados. Essa ênfase na construção de grupos comparáveis é o tema central da próxima seção.


A hipótese de independência condicional garante que, dentro de cada célula $X_i = x$, a seleção no tratamento é ``como se fosse aleatória''. Com isso, o efeito causal local é identificado diretamente pela diferença de médias observada:
\[
\delta(x)
=
\mathbb{E}\!\left[ Y_i \mid D_i = 1, X_i = x \right]
-
\mathbb{E}\!\left[ Y_i \mid D_i = 0, X_i = x \right]
=
\text{ATE}(x)
=
\text{ATT}(x).
\]

Uma vez identificados os efeitos causais locais $\delta(x)$, os parâmetros agregados surgem naturalmente como médias ponderadas desses efeitos. O ATE pode ser escrito como:
\[
\text{ATE}
=
\sum_x \text{ATE}(x)\,\mathbb{P}(X_i = x),
\]
ou seja, uma média ponderada dos efeitos locais, em que os pesos refletem a distribuição das características na população.

De forma análoga, o ATT é dado por:
\[
\text{ATT}
=
\sum_x \text{ATT}(x)\,\mathbb{P}(X_i = x \mid D_i = 1),
\]
isto é, a mesma média de efeitos locais, mas utilizando agora a distribuição das covariadas entre os tratados.


Se o vetor de covariadas $X_i$ fosse simples, a estimação de $\delta(x)$ e sua posterior agregação seriam diretas. Com apenas uma covariada, como sexo, por exemplo, bastaria calcular a diferença de médias condicional em ser mulher e a diferença de médias condicional em ser homem, e então agregar esses efeitos utilizando as proporções de mulheres e homens na amostra. Trata-se de um caso simples em que não seria necessário recorrer a técnicas de pareamento para estimar o efeito causal.

Na maior parte das aplicações empíricas, contudo, o vetor $X_i$ é substancialmente mais complexo. No caso do experimento do \textit{National Supported Work} (NSW), por exemplo, temos acesso a um conjunto de dez covariadas potencialmente relacionadas tanto à seleção no tratamento quanto aos resultados. Comparar indivíduos com perfis observáveis muito semelhantes nesse espaço de características é uma tarefa nada trivial.

O problema não é apenas a presença de ``muitos controles'', mas o fato de que o espaço de covariadas se torna rapidamente esparso à medida que sua dimensão aumenta. Nessas situações, encontrar tratados e controles verdadeiramente comparáveis em cada célula $X_i = x$ torna-se difícil ou mesmo impossível. É exatamente nesse ponto que o pareamento se torna útil: ao restringir explicitamente o conjunto de comparação, ele busca reconstruir comparabilidade local e evita que os efeitos causais sejam identificados por extrapolação entre unidades muito diferentes.


\section{Pareamento Exato}

O pareamento exato representa a forma mais direta e conceitualmente mais transparente de tornar tratados e controles comparáveis. A ideia é simples: para cada indivíduo tratado, buscamos no grupo de controle apenas indivíduos que tenham exatamente os mesmos valores em todas as covariadas utilizadas no pareamento. Quando isso é possível, a comparação deixa de depender de aproximações ou de noções graduais de proximidade. A comparabilidade é construída de forma literal.

Essa lógica implementa, da maneira mais fiel possível, a Hipótese de Independência Condicional apresentada anteriormente. Ao restringir a comparação apenas a indivíduos com o mesmo vetor de características observáveis $X_i$, eliminamos completamente as diferenças observáveis relevantes entre tratados e controles dentro de cada estrato. Com isso, a seleção no tratamento passa a ocorrer apenas dentro de grupos homogêneos em termos de $X_i$.

Para ver como isso funciona na prática, considere novamente um exemplo inspirado no contexto do NSW. Suponha que utilizemos apenas duas covariadas para o pareamento: escolaridade (baixa ou alta) e experiência prévia de trabalho (sim ou não). A Tabela 3.1 apresenta uma pequena amostra hipotética de indivíduos tratados e controles.

\begin{table}[h]
    \centering
    \begin{tabular}{lllll}
       Indivíduo  & Tratamento & Escolaridade & Experiência prévia & Salário \\\midrule
        A & 1 & Baixa & Não & 1.200\\
        B & 1 & Alta & Não & 1.400 \\
        C & 1 & Baixa & Sim & 1.300 \\
        1 & 0 & Baixa & Não & 900 \\
        2 & 0 & Alta & Não & 1.050\\
        3 & 0 & Baixa & Sim & 1.000 \\
        4 & 0 & Alta & Sim & 1.150\\
        \bottomrule
    \end{tabular}
    \caption{Exemplo numério pareamento exato}
    \label{tab:placeholder}
\end{table}

Aplicando o pareamento exato, cada tratado só pode ser comparado com controles que apresentem exatamente o mesmo perfil de escolaridade e experiência. O processo de emparelhamento gera então os pares apresentados na Tabela 3.2.

\begin{table}[h]
    \centering
    \begin{tabular}{lllll}
       Tratado  & Escolaridade & Experiência prévia & Controle correspondente & Salário controle \\\midrule
        A & Baixa &  Não  & 1 & 900\\
        B & Alta   & Não &  2 & 1.050 \\
        C & Baixa  & Sim  & 3 & 1.000 \\
        \bottomrule
    \end{tabular}
    \caption{Exemplo numério pareamento exato}
    \label{tab:placeholder}
\end{table}

O controle 4 não é utilizado, pois não existe nenhum tratado com escolaridade alta e experiência prévia simultaneamente. Da mesma forma, se algum tratado não encontrasse nenhum controle com perfil idêntico, ele simplesmente seria excluído da análise por falta de contrafactual observável adequado. Com base nesses emparelhamentos, os efeitos individuais são apresentados na Tabela 3.3.

\begin{table}[H]
    \centering
    \begin{tabular}{llll}
       Tratado  & Salário tratado & Salário controle & Efeito  \\\midrule
        A & 1.200 &   900    & 300 \\
        B & 1.400   & 1.050 &  350 \\
        C & 1.300  &  1.000  & 300  \\
        \bottomrule
    \end{tabular}
    \caption{Exemplo numério pareamento exato}
    \label{tab:placeholder}
\end{table}

O efeito médio sobre os tratados (ATT) é então simplesmente a média dessas diferenças:

$$
\widehat{ATT} = \frac{300 + 350 + 300}{3} = 316,7.
$$

Nesse ponto, fica claro o principal atrativo do pareamento exato: a interpretação do efeito é direta e livre de ambiguidades. Cada tratado é comparado com um controle que é idêntico em todas as características observáveis relevantes. O método não requer métricas de distância, funções de ponderação ou parâmetros de suavização. Toda a credibilidade da comparação vem da homogeneidade dentro dos estratos.

No entanto, essa mesma força é também a sua principal limitação. À medida que o número de covariadas aumenta, o número de combinações possíveis cresce rapidamente. Mesmo com poucas variáveis binárias, o número de estratos explode, e muitos deles passam a conter apenas tratados ou apenas controles. Nesses casos, o pareamento exato se torna inviável, pois grande parte da amostra precisa ser descartada por falta de pares exatos.

Esse problema é uma manifestação direta da \textbf{maldição da dimensionalidade}. Quanto mais refinado é o critério de igualdade exata entre tratados e controles, mais difícil se torna encontrar sobreposição suficiente entre os grupos. Na prática, o pareamento exato tende a funcionar bem apenas quando o número de covariadas é pequeno, quando elas são discretas e quando há forte sobreposição entre tratados e controles. Por essa razão, embora o pareamento exato seja conceitualmente o método mais limpo, ele raramente é utilizado isoladamente em aplicações empíricas com muitas covariadas.

\section{Pareamento como Médias Ponderadas}

Depois de vermos como o pareamento exato funciona e por que ele rapidamente se torna inviável à medida que o número de covariadas cresce, surge naturalmente a pergunta: o que fazer quando não existe ninguém no grupo de controle exatamente igual a um tratado? A resposta prática é simples e poderosa: em vez de utilizar um único controle, passamos a usar vários controles ``parecidos'', combinados por meio de uma média. É exatamente nesse ponto que o pareamento passa a ser entendido como um procedimento de médias ponderadas.

É importante enfatizar que essa transição para médias ponderadas não relaxa a hipótese de independência condicional. A CIA continua sendo assumida no mesmo conjunto de covariadas observáveis. O que se relaxa não é a identificação causal, mas apenas a exigência operacional de igualdade exata entre tratados e controles. Métodos de pareamento aproximado lidam com o problema de escassez de observações comparáveis no espaço de $X$, e não com viés decorrente de fatores não observáveis.

A ideia é a seguinte: para cada indivíduo tratado, queremos construir um valor de comparação que represente, da melhor forma possível, como ele teria se comportado sem o tratamento. Como não encontramos um “clone” perfeito no grupo de controle, usamos vários indivíduos semelhantes. Alguns deles são mais parecidos, outros menos. Os mais parecidos devem influenciar mais esse valor de comparação; os menos parecidos, menos. O contrafactual prático do tratado passa a ser, então, uma média ponderada dos resultados observados dos controles mais próximos.

No contexto do NSW, pense em um indivíduo tratado com certa idade, escolaridade, renda passada e histórico de desemprego. Não existe ninguém no grupo de controle com exatamente esse mesmo perfil. Ainda assim, encontramos vários controles com perfis próximos: alguns têm a mesma escolaridade, mas são um pouco mais velhos; outros têm a mesma idade, mas escolaridade ligeiramente diferente; outros têm renda passada parecida. O pareamento passa a combinar todos esses indivíduos para formar um “grupo sintético” de comparação. O salário contrafactual desse tratado é aproximado pela média dos salários desses controles, com pesos que refletem o grau de semelhança entre eles.

Essa lógica pode ser escrita de forma simples como

$$
\Tilde{Y}_i = \sum_{j \in \mathcal{C}} w_{ij} Y_j
$$

\noindent em que $\mathcal{C}$ é o conjunto de controles e os pesos $w_{ij}$ somam 1. A diferença entre os métodos de pareamento não está nessa expressão em si --- ela aparece na forma como esses pesos são definidos.

No pareamento exato, esses pesos são extremos: só recebem peso os controles que são exatamente iguais ao tratado, e todos eles recebem o mesmo peso. Nos métodos aproximados, que estudaremos a seguir, os pesos passam a variar de forma contínua. Controles muito parecidos recebem peso alto; controles pouco parecidos recebem peso baixo; controles muito diferentes recebem peso zero. Toda a prática do pareamento pode ser entendida, portanto, como diferentes regras para escolher esses pesos.

Essa interpretação como médias ponderadas ajuda a entender por que o pareamento é, ao mesmo tempo, flexível e transparente. Ele é flexível porque permite construir comparações mesmo quando não há correspondência exata. E é transparente porque deixa explícito quem está sendo usado para comparar com quem e com que peso cada observação entra no cálculo do contrafactual.

Ela também ajuda a esclarecer a relação entre pareamento e regressão. Na regressão, a forma como as observações são ponderadas na comparação entre tratados e controles é determinada implicitamente pela forma funcional do modelo. No pareamento, essa ponderação é construída explicitamente com base na similaridade entre indivíduos. Em vez de usar uma relação global entre $Y_i$, $D_i$ e $X_i$, o pareamento constrói comparações localizadas, indivíduo por indivíduo.

No exemplo do NSW, isso significa que um jovem tratado com baixa escolaridade e histórico recente de desemprego será comparado principalmente com controles com perfil semelhante, e não com a média geral do grupo de controle. O exercício de comparação passa a respeitar muito mais de perto a heterogeneidade dos indivíduos observados na amostra.

Ao mesmo tempo, essa forma de construir os contrafactuais envolve necessariamente um \textit{trade-off} entre precisão e estabilidade. Se usamos poucos controles muito próximos, a comparação tende a ser mais precisa, mas mais variável. Se usamos muitos controles, a comparação fica mais estável, mas corre o risco de misturar indivíduos pouco comparáveis. Esse equilíbrio entre “proximidade” e “quantidade de informação” estará presente em todos os métodos de pareamento que veremos a seguir.

Com essa interpretação em mente, estamos prontos para avançar para os critérios concretos usados para definir esses pesos na prática. A próxima etapa é formalizar o conceito de proximidade entre indivíduos, por meio de medidas de distância no espaço das covariadas, e estudar como essa noção dá origem aos principais métodos operacionais de pareamento, como o pareamento por vizinho mais próximo (\textit{Nearest Neighbor}) e por \textit{kernel}. Em seguida, veremos como o \textit{propensity score} surge como uma forma particularmente poderosa de organizar essa ponderação quando o número de covariadas é grande.

